% SHARD-ID: RC-04W-GENUS-TWO-BOND
% SHARD-TITLE: The Phantasm XIII: the Class-Group Bond at Genus Two. The Curve y^2 = x^5 + x^3 + x^2 - 2 over F_5, Its Fifteen Classes and Fourteen L-Functions, the Pointed Class-Field Cover of Genus 16, the Two-Parameter Metric Cone, the CM Field and Its Reference Polarisation, No Rational Symmetric Marking, and the Degree Cut That Is Not H^1
% SHARD-SUMMARY: The genus-two curve of shard 06f has counts (3, 31, 117, 619), P(T) = 1 - 3T + 7T^2 - 15T^3 + 25T^4, Pic^0(F_5) = Z/15 generated by [(1,1) - infinity], theta classes {0, 1, 14}, special counts (1, 3, 20); all fourteen nontrivial class characters have L(T, chi_k) = 1 + (1 + zeta^k + zeta^{-k}) T + 5 T^2 with zeros of modulus 5^{-1/2} exactly, functional equation L_k(T) = 5 T^2 L_{15-k}(1/(5T)) (constant chi(kappa) = 1), and the pointed Hilbert class-field cover has genus 16 with P_Y = P prod L(T, chi): the bond's Fourier decomposition labels the character summands of H^1(Y), without an isomorphism of spaces.
% SHARD-SUMMARY: The trivial Frobenius block carries a two-parameter cone of invariant positive metrics (one ratio modulo scale, the first non-vacuous instance); deck-invariant metrics on the seven inverse character pairs are diag(a_k, b_k, b_k, a_k). The CM field is K = Q(sqrt 21, delta), delta^2 = (-25 + 3 sqrt 21)/2, O_K = Z[pi, 5/pi] maximal of discriminant 21^2 * 109; the principal Riemann form Tr(xi x conj y) gives weights |phi_j(xi)| with reference ratio (25 + 3 sqrt 21)/(2 sqrt 109); equal weights at the two CM places are impossible for any rational marking (they would force an imaginary quadratic subfield), and the unit (5 + sqrt 21)/2 rescales the ratio by (527 + 115 sqrt 21)/2, so the bond-to-H^1 comparison must be supplied before a symmetric-ray statement has meaning.
% SHARD-SUMMARY: The degree shift compressed to the special degrees is a multiplicity-15 coisometry (all degrees n <= 2) or a 45-dimensional nilpotent with characteristic function z^3 I_15 (degrees 0, 1, 2): it does not carry P; the theta-cyclic cut has a rank-two exit; the scalar Blaschke model K_{P/f_F} has the Cauchy Gram, outside the invariant cone (normality defect 1.3238593874), and any d-mode invariant-metric channel has defect rank d, incompatible with one exit. The two ends of the GL_1 bond are not two cusps; the genus-two H-CLASS fold is stated as a conditional 2 x 2-block prediction.
% SHARD-KEYWORDS: genus two, class group, Mumford representation, Cantor arithmetic, theta divisor, class characters, Hecke L-function, Hilbert class field, Artin reciprocity, Riemann--Hurwitz, CM field, Riemann form, polarisation, unit group, metric cone, symmetric ray, degree shift, coisometry, nilpotent, Blaschke product, Cauchy Gram
% SHARD-NOTES: none
% SHARD-SCRIPTS: scripts/genus_two_bond.py

\section{The Phantasm XIII: the Class-Group Bond at Genus Two}
\label{sec:genus-two-bond}

Round of 2026-09-22 (brief \texttt{notes/genus-two-bond/astra-brief.md}, prover \texttt{notes/genus-two-bond/astra-proofs.md},
numerics \texttt{scripts/genus\_two\_bond.py}, review \texttt{notes/reviews/genus-two-bond-2026-09-22.md}). The first
non-vacuous test of the metric question left open by \cref{obs:gl1-bond-no-metric}, on the curve $C: y^2 = x^5+x^3+x^2-2$
over $\mathbb F_5$ of \cref{sec:ring-norm-genus-two}, with $P_0 = \infty$, $P_\pm = (1,\pm1)$ (i.e. $(1,1)$ and $(1,4)$),
$G = [P_+-P_0]$, $\zeta = e^{2\pi i/15}$, $A_k = 1+\zeta^k+\zeta^{-k}$, and $\Omega_\theta$ the theta bond state. % bare-ok % bare-ok

\subsection{The finite data and the fourteen $L$-functions}

\begin{proposition}[The class-group bond of the genus-two curve]
\label{prop:genus-two-class-bond-data}
\claimstatus{prop:genus-two-class-bond-data}{proved}
$N_k = 3,31,117,619$ (with the one rational point over $\infty$ of the odd-degree model), so $P(T) = 1-3T+7T^2-15T^3+25T^4$,
$f_F(z) = (z^2-b_1z+5)(z^2-b_2z+5)$ with $b_{1,2} = (3\pm\sqrt{21})/2$, four distinct roots of modulus exactly $\sqrt5$
($b_j^2<20$; independent of \cref{asm:weil-curves}), $h = P(1) = 15$. $\mathrm{Pic}^0(F_5) = \mathbb Z/15$ generated by $G$ (Mumford
representation, Cantor addition audited by hand: fifteen reduced pairs, $[P_--P_0] = 14G$); the effective classes are
$\{0\}$ in degree $0$, the theta set $\{0,1,14\}$ in degree $1$ (each with $h^0 = 1$), and all fifteen in degree $2$ with
$h^0 = 2$ for the principal class (six representatives: $2P_0$ and the five fibres of $x$) and $1$ otherwise; the counts are
$a_0,a_1,a_2 = 1,3,20$ ($6+14$ closed points of degree two) and $a_n = 15(5^{n-1}-1)/4$ for $n\ge3$, and
$(1-6T+5T^2)\sum a_nT^n = P(T)$ exactly: the instance of \cref{prop:class-group-bond}. Verified: \cref{num:genus-two-bond}.
\end{proposition}

\begin{theorem}[The pointed class-field cover and the character decomposition of its $H^1$]
\label{thm:pointed-class-field-cohomology}
\claimstatus{thm:pointed-class-field-cohomology}{proved}
For $X/\mathbb F_q$ smooth projective geometrically connected of genus $g\ge1$ with a rational point $O$, $A = \mathrm{Pic}^0(X)(\mathbb F_q)$,
$h = |A|$, and characters extended by $\chi(O) = 1$: the maximal abelian extension unramified everywhere in which $O$ splits
completely is geometric of degree $h$ with Galois group $A$ (Artin reciprocity for $\mathrm{Pic}(X)/\mathbb Z[O]\cong A$;
\texttt{1012.3223:main.tex:318-332,1235-1244}, \texttt{1112.5826:main.tex:534}), of genus $g_Y = 1+h(g-1)$ (Riemann--Hurwitz,
\'etale);
$L_X(T,\chi) = \sum_n\sum_{c\in\mathrm{Pic}^n}\chi(c-nO)(q^{h^0(c)}-1)/(q-1)T^n$ is, for $\chi\ne1$, a polynomial of degree $2g-2$ with
constant term $1$, leading coefficient $q^{g-1}\chi(\kappa)$, $\kappa = [K_X-(2g-2)O]$, and functional equation
$L_X(T,\chi) = \chi(\kappa)q^{g-1}T^{2g-2}L_X(1/(qT),\chi^{-1})$ (no extra sign; Riemann--Roch and $D\mapsto K_X-D$); with $\ell\ne p$ and
$\mathcal L_\chi$ the rank-one summands of $p_{*}\bar{\mathbb Q}_\ell$ (geometric Frobenius eigenvalue $\chi([Q-\deg Q\,O])$),
$H^1(Y) = \bigoplus_\chi V_\chi$, $V_1 = H^1(X)$, $\dim V_\chi = 2g-2$, $\det(1-TF|V_\chi) = L_X(T,\chi)$, and
$P_Y = P_X\prod_{\chi\ne1}L_X(T,\chi)$ (Grothendieck's trace formula, not byte-cited). For $C$: $\kappa = 0$ ($K_C\sim2P_0$),
$L_k(T) = 1+A_kT+5T^2$ for $k = 1,\dots,14$ ($A_k\in\{(1\pm\sqrt5)/2, 0, 1+\zeta^k+\zeta^{-k}\}$ by order $5, 3, 15$), zeros
$(-A_k\pm i\sqrt{20-A_k^2})/10$ of modulus $5^{-1/2}$ exactly ($A_k^2\le9<20$), $L_k(T) = 5T^2L_{15-k}(1/(5T))$, and
$P_Y = P\,(1+5T^2)^2(1+T+9T^2+5T^3+25T^4)^2(1+5T+25T^2+70T^3+195T^4+350T^5+625T^6+625T^7+625T^8)^2$ of degree $32 = 2g_Y$,
$g_Y = 16$, $N_1(Y) = 15$ ($\infty$ splits completely, $P_\pm$ do not). The predecessor's ``a single $L$-function is the
numerator of the cover divided by that of $C$'' (\texttt{notes/gl1-bond/proofs.md:99-101}) must read ``the product over all
nontrivial characters''.
\end{theorem}

\begin{observation}[The bond labels $H^1(Y)$; it is not $H^1(Y)$]
\label{obs:bond-h1-labels}
\claimstatus{obs:bond-h1-labels}{proved}
Fourier transformation on $\ell^2(A)$ supplies the labels $\chi$ and, applied to the divisor-count sequences, the
$L$-functions, whose cohomological realisation is $H^1(Y)$ with the multiplicities above. There is no isomorphism
$\ell^2(A)\cong H^1(Y)$ ($15$ against $32$; the degree bond is infinite-dimensional) and the degree shift is not Frobenius
(\cref{thm:degree-cut-not-h1}).
\end{observation}

\subsection{The Frobenius channel and the metric cone}

\begin{proposition}[The genus-two metric cone: two weights on the trivial block; inverse pairs on the cover]
\label{prop:genus-two-metric-cone}
\claimstatus{prop:genus-two-metric-cone}{proved}
A finite-dimensional operator is unitary in some positive metric iff it is semisimple with unimodular spectrum. On the
trivial block with $D = \mathrm{diag}(\alpha_{j,\pm}/\sqrt5)$, $(D^{*}HD-H)_{ab} = (\bar\mu_a\mu_b-1)H_{ab}$ vanishes exactly on the
diagonal (distinct unimodular roots), so the invariant positive metrics are $\mathrm{diag}(h_{1,+},h_{1,-},h_{2,+},h_{2,-})$, and the
\emph{additional} reality condition $H(cx,cy) = \overline{H(x,y)}$ for the antilinear $c$ exchanging $\alpha_{j,\pm}$ gives
$\mathrm{diag}(h_1,h_1,h_2,h_2)$: two parameters, one ratio modulo scale, with no genus-two symmetry exchanging $j = 1,2$
(the D2 single ray of \cref{thm:arithmetic-metric} used its transitive parity on the normalised cavity vectors).
``$F$ is real on $H^1(C,\mathbb Q_\ell)$'' is not the reason: a $\mathbb Q_\ell$-space has no distinguished conjugation; the real
structure lives on the polarised characteristic-zero lift or on a real realisation of the joint deck/Frobenius
representation. On the cover, conjugation sends $V_{\chi_k,\lambda_{k,+}}$ to $V_{\chi_{15-k},\lambda_{k,-}}$; for deck-invariant
metrics distinct characters are orthogonal and each inverse pair $(k,15-k)$ carries $\mathrm{diag}(a_k,b_k,b_k,a_k)$;
$F$-invariance alone does not force character orthogonality, since $L_k = L_{15-k}$ makes the common eigenspaces
two-dimensional.
\end{proposition}

\begin{proposition}[Two Frobenius-channel conventions: $32$ direct modes or $64$ square-root modes]
\label{prop:genus-two-frobenius-channel}
\claimstatus{prop:genus-two-frobenius-channel}{proved-conditional}
With $U = F/\sqrt5$ on the $32$-dimensional Frobenius module of the cover (one $4\times4$ and fourteen $2\times2$ blocks), $r = 5^{-1/4}$,
$Z' = rU$, $J' = \sqrt{1-5^{-1/2}}$, the renewal channel $\mathcal E(X) = 5^{-1/2}UXU^{*}+(1-5^{-1/2})\Tr(X)\reb$ has all relaxation
eigenvalues of modulus $5^{-1/2}$ and modal factors $(z-r\mu)/(1-r\bar\mu z)$, a direct application of
\cref{thm:frobenius-channel-expander}; the D2 cavity convention there ($U^2\simeq(F/\sqrt q)\otimes1_2$) instead needs $64$
modes $\mu = \pm\sqrt{\alpha/\sqrt5}$, eight for the trivial character and four per nontrivial one. Both are unitary after
choosing an invariant metric; neither selects one.
\end{proposition}

\begin{theorem}[The CM field, its maximal order, and the reference polarisation weights]
\label{thm:genus-two-cm-reference}
\claimstatus{thm:genus-two-cm-reference}{proved-conditional}
With $\beta = \pi+5/\pi$, $\delta = \pi-5/\pi$: $\beta^2-3\beta-3 = 0$, $K^+ = \mathbb Q(\sqrt{21})$ ($\sqrt{21} = 2\beta-3$),
$\delta^2 = 3\beta-17 = (-25+3\sqrt{21})/2$ with both conjugates negative, $K = K^+(\delta)$ a quartic CM field, $f_F$ irreducible
(mod $2$), the Jacobian simple and ordinary (\texttt{1701.07742:main.tex:551-555}); $O_{K^+} = \mathbb Z[\beta]$,
$O_K = \mathbb Z[\beta,\pi] = \mathbb Z[\pi,5/\pi]$ with basis $(1,\beta,\pi,\beta\pi)$, trace-matrix determinant $48069 = 21^2\cdot109$,
$[O_K:\mathbb Z[\pi]] = 5$ (relative discriminant $(3\beta-17)$ of norm $109$, squarefree); $K$ has no imaginary quadratic
subfield ($N_{K^+/\mathbb Q}(\delta^2) = 109$ is not a square) and is not normal (Galois group $D_4$). Under
\cref{asm:cm-hodge-type,asm:cm-polarised-lift}, after a rational $K$-module marking the lattice is a fractional ideal $I$, the
principal Riemann form is $E(x,y) = \Tr_{K/\mathbb Q}(\xi x\bar y)$ with $\bar\xi = -\xi$, $\xi I\bar I = D_K^{-1}$
(Shimura--Taniyama; Milne, Prop.\ 2.9; not byte-cited), and the polarisation-transported metric on $K\otimes\mathbb R$ is
$2\sum_j\im\phi_j(\xi)\re(x_j\bar y_j)$: weights $c_j = |\phi_j(\xi)|$, orthogonal blocks, Rosati $=$ conjugation, $\pi\bar\pi = 5$
unitarity (the dual period-form normalisation has weights $1/|\phi_j(\xi)|$). The reference lattice $I_0 = O_K$,
$\xi_0 = 1/(\sqrt{21}\,\delta)$ (the different is $(\sqrt{21}\delta)$) has integer unimodular $E_0$ in the basis above,
CM type $\Phi_0 = \{\alpha_{1,-},\alpha_{2,+}\}$, weights $c^0_j = 1/(\sqrt{21}d_j)$ ($0.0919994481$, $0.0495772273$) and ratio
$R_0 = d_2/d_1 = (25+3\sqrt{21})/(2\sqrt{109}) = 1.8556795747$, $R_0^2 = (407+75\sqrt{21})/218\in K^+$, $R_0\in\mathbb Q(\sqrt{21},\sqrt{109})\setminus K$.
\end{theorem}

\begin{proposition}[No rational marking gives equal weights; the unit rescaling]
\label{prop:no-rational-symmetric-marking}
\claimstatus{prop:no-rational-symmetric-marking}{proved-conditional}
For any nonzero purely imaginary $\xi = a\delta$, $a\in K^{+\times}$, the transported ratio is $|\sigma_1(a)/\sigma_2(a)|\,d_1/d_2$ and its
square lies in $K^+$; equal weights $|\phi_1\xi| = |\phi_2\xi|$ would make the two real conjugates of $\xi^2$ agree, so
$\xi^2\in\mathbb Q_{<0}$ and $K$ would contain an imaginary quadratic field, which it does not. Moreover the unit
$\varepsilon = (5+\sqrt{21})/2\in O_{K^+}^\times$ (norm $1$, both embeddings positive, $\varepsilon I = I$) rescales the ratio by
$\varepsilon_1^4 = (527+115\sqrt{21})/2$ without changing the lattice: infinitely many presentations before any real rescaling, so
there are not ``two candidate answers'' for an ideal class to select. The theta normalisation on the bond is computable
(Fourier components $19,27,95$ over $\sqrt{15}$ for $k = 0$ and $4,4A_k,20$ over $\sqrt{15}$ for $k\ne0$; total squared norm
$1101$; degree polynomials $15^{-1/2}(1+3T+20T^2)$ and $15^{-1/2}L_k(T)$) but gives one three-vector per character and no
components along the four Frobenius eigenlines; the $L$-values $L_k(1) = 6+A_k$, $L_k(5^{-1/2}) = 2+A_k/\sqrt5$ normalise
whole blocks, not eigenlines within a block.
\end{proposition}

\begin{observation}[The bond-to-$H^1$ comparison is the missing datum]
\label{obs:bond-h1-comparison-open}
\claimstatus{obs:bond-h1-comparison-open}{open}
``The symmetric ray of \cref{thm:arithmetic-metric} is the Hodge metric'' has no specified meaning for this curve until a
map, or a normalisation rule on geometric $H^1$, invariantly determined by the bond and compatible with $F$, is supplied;
neither the theta coefficients nor the $L$-values supply it. The finite arithmetic task for a marked lift is explicit
(ideal classes of $O_K$, Minkowski bound $<34$; solve $\xi I\bar I = D_K^{-1}$ modulo $(I,\xi)\sim(aI,\xi/(a\bar a))$; the
ordinary CM type by a $5$-adic condition), and would still need the comparison.
\end{observation}

\subsection{The degree cut is not $H^1$}

\begin{theorem}[The special-degree cut of the degree shift does not realise $H^1$; the correct scalar-exit model has a non-normal Gram]
\label{thm:degree-cut-not-h1}
\claimstatus{thm:degree-cut-not-h1}{proved}
The degree shift $S|j,n\rangle = |j,n+1\rangle$ on $\ell^2(A)\otimes\ell^2(\mathbb Z)$ involves neither $h^0$, $P$ nor $\theta$. Compressed to
all degrees $n\le2$ it is a multiplicity-$15$ backward unilateral shift (a coisometry, infinite-dimensional, no finite
characteristic determinant); compressed to degrees $0,1,2$ it is $I_{15}\otimes N_3$, nilpotent, $45$-dimensional, with
defect rank $15$ at each end, Gram $I_{45}$, characteristic function $z^3I_{15}$, $\det\Theta = z^{45}$: no nilpotent
restriction or quotient carries the nonzero Frobenius spectrum, and $N^{*}HN = r^2H$ forces $H = 0$. The theta-cyclic cut
$(\theta,T\theta,T^2\theta)$ has Gram $\begin{psmallmatrix}1101&282&195\\282&126&47\\195&47&39\end{psmallmatrix}$ ($\det = 254679$) and exit
defect of rank two ($\theta_2 = 5\theta_0$, $\theta_1$ independent): two exit coordinates, not one. The Perron tail is a
statement about the scalar count series (poles at $1$ and $1/5$, $P$ its transmission numerator), not a subspace whose
removal turns $S$ into $F$. The correct scalar-exit model built \emph{after} extracting $P$ is $K_B$, $B = P(w)/f_F(w)$ a
degree-four Blaschke product with zeros $\alpha_{j,\pm}/5$, whose normalised eigenvectors have the Cauchy Gram
$(4/5)/(1-\bar\rho_a\rho_b)$ with all off-diagonal entries nonzero (the identity of \cref{thm:h-exit-model}): outside the cone of
\cref{prop:genus-two-metric-cone}, normality defect $\max|(F/\sqrt5)^{*}G(F/\sqrt5)-G| = 1.3238593874$; the square-root
convention $B(z^2)$ has the same obstruction. Structurally, a $d$-mode invariant-metric channel $Z' = rU$ has defect
$(1-r^2)I_d$ of rank $d$, while a scalar exit has defect rank $\le1$: incompatible for $d = 4$ or $8$, the genus-two form of
the D2 obstruction of \cref{sec:elliptic-cavity-channel}.
\end{theorem}

\begin{observation}[Two ends, one cusp; the genus-two H-CLASS fold]
\label{obs:two-ends-fold}
\claimstatus{obs:two-ends-fold}{sketched}
\Cref{prop:siegel-constant-term} is the theorem: the two theta terms of the level-one constant term are exchanged by
$F(t)\mapsto t^{-1}F(1/t)$ at fixed $y$. $\mathbb R_+^\times$ has two topological ends, not two cusps (\cref{sec:gl1-bond} says one
cusp); ``the automorphism folds the two GL$_1$ cusps into the GL$_2$ cusp'' is a metaphor, open as an identification. On
this bond the Riemann--Roch involution is $(j,n)\mapsto(-j,2-n)$ with $\theta_n(j) = 5^{n-1}\theta_{2-n}(-j)$, sending $k$ to $15-k$;
there are fifteen cusp classes for $\mathrm{GL}_2(O(C\setminus P_0))$ (\cref{cit:lorscheid-cusp-count}), and if the scattering matrix
has the H-CLASS group-matrix form its Fourier transform is one scalar block and seven blocks
$\begin{psmallmatrix}0&m_k\\m_{15-k}&0\end{psmallmatrix}$ with eigenvalues $\pm\sqrt{m_km_{15-k}}$, where H-CLASS predicts $L$-ratios:
this curve supplies fourteen nonconstant quadratic $L$-factors where genus one supplied constants. Proving that this
quotient has that scattering matrix is open.
\end{observation}

\begin{numerical}[Numerics lane, the genus-two bond]
\label{num:genus-two-bond}
\claimstatus{num:genus-two-bond}{numerical}
\texttt{scripts/genus\_two\_bond.py} ($72$ assertions, $9$ s, sympy/numpy, exact where the claim is exact;
\texttt{outputs/genus\_two\_bond.txt}; ledger in \texttt{notes/genus-two-bond/numerics.md}; $72$ pass; written without reading the
prover's scripts): the counts with own $\mathbb F_{5^k}$ arithmetic; $P$, $h$, the roots and their moduli; Cantor arithmetic from
scratch (fifteen reduced pairs, $G$ of order $15$, all $225$ additions); the effective classes and $(1,3,20)$, the $Z(T)$ series;
all fourteen $L_k$ in $\mathbb Q(\zeta_{15})$ with their zeros, functional equations and the exact degree-$32$ product; the CM
field data ($\beta,\delta$, the trace matrix with determinant $48069$, $\mathrm{disc}\,\mathbb Z[\pi] = 25\cdot48069$, $N(\delta^2) = 109$,
ordinarity); $E_0$, the weights, the ratio and its square, the unit identity; the theta Fourier weights and total $1101$;
the theta-cyclic Gram, its determinant and defect rank; the nilpotent cut; the Cauchy Gram with the normality defect
$1.3238593874$; the cone counts ($60/32/16$ real parameters on $H^1(Y)$). The review adds the $D_4$ Galois group, the
index-$5$ generator $5/\pi+\pi^2-\pi+3$, and the $\varepsilon$ identity.
\end{numerical}

\begin{observation}[What this changes; next]
\label{obs:genus-two-next}
\claimstatus{obs:genus-two-next}{sketched}
Genus two makes the metric question real (a ratio modulo scale) and shows what it is not: neither the bond's theta
normalisation nor the $L$-values nor a rational marking fixes it, and the class-field cover's $H^1$ is labelled, not
carried, by the bond. The same lesson as \cref{sec:deninger-lps}: positivity is supplied by a polarised lift (Petersson
or the Riemann form), the bond supplies labels and $L$-functions. Next: the genus-two quotient graph and its scattering
matrix (the conditional H-CLASS fold), and the comparison map, if any, from the bond to $H^1$.
\end{observation}
